\CatchFileBetweenTags{\AlphaInvCodataExperimentalVal}{constants.tex}{AlphaInvCodataExperimentalVal}
\CatchFileBetweenTags{\AlphaInvCodataSigma}{constants.tex}{AlphaInvCodataSigma}
\CatchFileBetweenTags{\AlphaInvMorelExperimentalVal}{constants.tex}{AlphaInvMorelExperimentalVal}
\CatchFileBetweenTags{\AlphaInvMorelSigma}{constants.tex}{AlphaInvMorelSigma}
\CatchFileBetweenTags{\AlphaInvParkerExperimentalVal}{constants.tex}{AlphaInvParkerExperimentalVal}
\CatchFileBetweenTags{\AlphaInvParkerSigma}{constants.tex}{AlphaInvParkerSigma}
\CatchFileBetweenTags{\AlphaInvFanExperimentalVal}{constants.tex}{AlphaInvFanExperimentalVal}
\CatchFileBetweenTags{\AlphaInvFanSigma}{constants.tex}{AlphaInvFanSigma}

\section{Empirical Correspondences}
\label{sec:empirical_correspondences}
$Z_{\rm geo}$ is a dimensionless property of the substrate. To compare it with experiment we must say which observable it is. This section is outside the main theorem: it adds one physical hypothesis and follows its consequences.

\subsection{Electromagnetic Matching}
\label{sec:electromagnetic_matching}
Suppose coarse graining the substrate yields a local $U(1)$ gauge theory. In rationalized natural units, with the connection normalized so that a unit charge enters as $D_\mu=\partial_\mu+iA_\mu$, write the leading term as $-\frac{1}{4g^2}F_{\mu\nu}F^{\mu\nu}$; then $\alpha=g^2/4\pi$\cite{tong_qft_2007}.

\textbf{Matching hypothesis.} The substrate's net impedance is the zero-momentum Maxwell coefficient:
\begin{equation}
\label{eq:maxwell_matching}
\frac{4\pi}{g^2(0)}=Z_{\rm geo},\qquad\text{so}\qquad \alpha^{-1}(0)=Z_{\rm geo},\quad \alpha(0)=\kappa.
\end{equation}
The quiescent baseline has no external flux and no running fields, so it corresponds to the infrared limit $Q^2\to0$. The algebra and the inherited bounds are checked in Lean (\texttt{electromagnetic\_matching}, \texttt{physical\_inverse\_coupling\_bounds}). The coefficient $4\pi$ is part of the hypothesis: had the calculation produced $Z_{\rm geo}=C/g^2(0)$, the prediction would be $\alpha^{-1}(0)=(4\pi/C)Z_{\rm geo}$ (\texttt{matching\_coefficient}). Deriving \Cref{eq:maxwell_matching} requires the substrate's coupling to a background $U(1)$ connection and its quadratic response, which this paper does not compute.

\subsection{The Fine-Structure Constant}
\label{sec:fine_structure}
\begin{table}[t]
\centering
\resizebox{\columnwidth}{!}{%
\begin{tabular}{lcc}
\toprule
\textbf{Source} & \textbf{$\alpha^{-1}(0)$} & \textbf{$Z_{\rm geo}-\alpha^{-1}$} \\
\midrule
Morel (2020), Rb recoil & \AlphaInvMorelExperimentalVal & $\AlphaInvMorelSigma\,\sigma$ \\
Fan (2023), electron $g-2$ & \AlphaInvFanExperimentalVal & $\AlphaInvFanSigma\,\sigma$\\
Parker (2018), Cs recoil & \AlphaInvParkerExperimentalVal & $\AlphaInvParkerSigma\,\sigma$ \\
CODATA (2022) & \AlphaInvCodataExperimentalVal & $\AlphaInvCodataSigma\,\sigma$ \\
\bottomrule
\end{tabular}%
}
\caption{$Z_{\rm geo}$ against determinations of $\alpha^{-1}(0)$. Deviations use the certified value of $Z_{\rm geo}$ and each source's quoted uncertainty (Lean: \texttt{alpha\_deviations}).}
\label{tab:alpha_comparison}
\end{table}
Under the matching hypothesis the index predicts $\alpha^{-1}(0)$ with no adjustable parameter (\Cref{tab:alpha_comparison}). It agrees with the rubidium recoil measurement of Morel et al.\ within $0.58\,\sigma$ and with CODATA 2022 within $1.68\,\sigma$. The determinations themselves disagree: the Parker cesium recoil and Morel rubidium recoil values differ by more than $5\,\sigma$, and CODATA's uncertainty reflects that tension.

\paragraph{Quantization Gap hypothesis.} The substrate has finite channel capacity, while perturbative QED sums ever larger sets of diagrams. We conjecture that high-loop extractions of $\alpha$, such as from the electron $g-2$, deviate systematically from the kinematic baseline because the continuum expansion eventually exceeds the substrate's capacity. The Fan $g-2$ value lies $3.1\,\sigma$ below the index while Morel agrees within $0.6\,\sigma$, which is the pattern the hypothesis expects. It is not yet a prediction: that requires a mapping from diagrams to channel states, an onset loop order, and the sign and size of the correction to a named observable. Until then the table tests the matching hypothesis, not the gap.

\paragraph{Falsification.} A robust determination of $\alpha^{-1}(0)$ several standard deviations from $Z_{\rm geo}$, with its theoretical corrections under control, would refute the conjunction of the postulates and the matching hypothesis.

\subsection{Elementary Charge, Resistance, and Field Scales}
\label{sec:elementary_charge}
\label{sec:channel_resistance}
\label{sec:load_limit}
The matching hypothesis carries over to the observer's unit system through standard identities. With the Planck charge $q_P=\sqrt{4\pi\epsilon_0\hbar c}$,
\begin{equation}
e=\frac{q_P}{\sqrt{Z_{\rm geo}}},\qquad \frac{e}{q_P}=\sqrt{\kappa}.
\end{equation}
With the impedance of free space $Z_0=1/(\epsilon_0c)$, the von Klitzing constant is
\begin{equation}
R_K=\frac{h}{e^2}=\frac{Z_0}{2}\,Z_{\rm geo}
\end{equation}
(\texttt{resistance\_identity}): the quantum of Hall resistance is half the impedance of free space times the substrate's net impedance per cycle. With the reduced Compton wavelength $\bar\lambda=\hbar/m_ec$ and the Schwinger field $E_c=m_e^2c^3/e\hbar$\cite{schwinger_gauge_1951},
\begin{equation}
\frac{e}{4\pi\epsilon_0\bar\lambda^2}=\frac{E_c}{Z_{\rm geo}}
\end{equation}
(\texttt{schwinger\_scale\_identity}). These follow from the one hypothesis $\alpha^{-1}(0)=Z_{\rm geo}$ and are not independent confirmations of it. In the 2019 SI, $e$ and $h$ are exact, so they constrain the relation between $\alpha$ and the measured $\epsilon_0$ (equivalently $\mu_0$)\cite{mohr_codata_2025}. Charge quantization, the Hall plateaus\cite{thouless_quantized_1982,kohmoto_topological_1985}, and pair-production dynamics need further structure.
